\begin{lemma}\label{minremiscusp}
Let \(\xi\) be a minimal \(\ssee\)-removable ribbon in a skew shape \(\tau\), with \(\cont(\xi) \in \Psi\). Then \(\xi\) is cuspidal.
\end{lemma}
\begin{proof}
Take any \(\ssee\)-removable ribbon \(\nu \subsetneq \xi\) in \(\xi\). Since \(\xi\) is \(\ssee\)-removable in \(\tau\), we have that \(\nu\) is \(\ssee\)-removable in \(\tau\). By the minimality of \(\xi\), we have \(\cont(\nu) \succeq \cont(\xi)\), and by Lemma~\ref{notapprox}\eqref{cond:notapprox2} we have \(\cont(\nu) \not \approx \cont(\xi)\), so \(\cont(\nu) \succ \cont(\xi)\). Thus \(\xi\) is cuspidal by Lemma~\ref{CuspForSkewHook}.
\end{proof}

\begin{lemma}\label{mdeltaribbon}
If \(\nu\) is a ribbon of content \(m\delta\), then \(\nu\) has an \(\ssee\)-removable ribbon of content \(\delta\).
\end{lemma}
\begin{proof}
We go by induction on \(m\). The base case \(m=1\) is trivial, so assume \(m>1\) and make the induction assumption on \(m' < m\). By Lemma~\ref{hookispath} there exists a \(\no/\ea\) path \((z^i)_{i=1}^{me}\) such that \(\nu = \{z^i\}_{i=1}^{me}\). If \(z^{e+1} = \no z^e\), it follows by Lemma~\ref{xiuv} that \(\xi^\nu_{z^1, z^e}\) is \(\ssee\)-removable in \(\nu\), and we have \(\cont(\xi^\nu_{z^1,z^e}) = \delta\) by Lemma~\ref{hookcont}. On the other hand, if \(z^{e+1} = \ea z^e\), it follows by Lemma~\ref{xiuv} that \(\xi^\nu_{z^{e+1}, z^{me}}\) is \(\ssee\)-removable in \(\nu\), and we have \(\cont(\xi^\nu_{z^{e+1}, z^{me}}) = (m-1)\delta\) by Lemma~\ref{hookcont}. Then by induction there exists an \(\ssee\)-removable ribbon \(\omega\) in \(\xi^\nu_{z^{e+1}, z^{me}}\) of content \(\delta\). As \(\omega\) is \(\ssee\)-removable in \(\xi^\nu_{z^{e+1}, z^{me}}\), it is \(\ssee\)-removable in \(\nu\). This completes the induction step, and the proof.
\end{proof}


\begin{lemma}\label{minremexist}
If \(\tau\) is a skew shape, then \(\tau\) has a minimal \(\ssee\)-removable ribbon \(\xi\) such that \(\cont(\xi) \in \Psi\).
\end{lemma}
\begin{proof}
Let \(\nu\) be a minimal \(\ssee\)-removable ribbon in \(\tau\). Then \(\cont(\nu) \in \Phi_+\) by Lemma~\ref{hookcont}. If \(\cont(\nu) \in \Phi_+^\re\), we are done by (\ref{PsiDef}), so assume otherwise. Then 
by (\ref{allimag}), we have \(\cont(\nu) = m\delta\) for some \(m\in \NN\).
Then by Lemma~\ref{mdeltaribbon}, \(\nu\) has an \(\ssee\)-removable ribbon \(\xi\) of content \(\delta \in \Psi\). But then \(\xi\) is \(\ssee\)-removable in \(\tau\), and \(\delta \approx m\delta\), so \(\xi\) is minimal in \(\tau\) as well.
\end{proof}

The next key lemma establishes that consecutively \(\ssee\)-removed minimal ribbons are non-decreasing in the convex preorder. 

 \begin{lemma}\label{remsinorder}
Let \(\tau\) be a skew shape. Let \(\xi_1\) be a minimal \(\ssee\)-removable ribbon in \(\tau\), and let \(\xi_2\) be a minimal \(\ssee\)-removable ribbon in \(\tau'=\tau\backslash \xi_1\). Then \(\cont(\xi_2) \succeq \cont(\xi_1)\).
 \end{lemma}
 \begin{proof}
First, note that if \(\xi_1 \sqcup \xi_2\) is disconnected, then \(\xi_2\) is an \(\ssee\)-removable ribbon in \(\tau\), so by minimality of \(\xi_1\), we would have \(\cont(\xi_2) \succeq \cont(\xi_1)\), as desired. Thus we may assume \(\xi_1 \sqcup \xi_2\) is a connected skew shape, and therefore \(\xi_1, \xi_2\) belong to the same connected component of \(\tau\). 
  We have \(\xi_1 = \xi^\tau_{u,v}\) for some \((u,v) \in \textup{Rem}_\tau\), and \(\xi_2 = \xi^{\tau'}_{z,w}\) for some \((z,w) \in \textup{Rem}_{\tau'}\), where \(u,v,z,w\) belong to the same connected component of \(\tau\).  
We consider the six possible cases of arrangements of the nodes \(u,v,z,w\) separately.
 
 {\em Case 1: Assume \(u \nearrow v \nearrow z \nearrow w\).} First, note that since \(\xi_1 \sqcup \xi_2\) is connected, we must have \(z \in \{\no v, \ea v\}\). But, since \((u,v) \in \textup{Rem}_\tau\), we have \(\ea v \notin \tau\), so \(z=\no v\). 
By Lemma~\ref{remhookswap}\eqref{cond:remhookswap1} we have that \((u,w) \in \textup{Rem}_\tau\) and \(\xi^\tau_{u,w}\) is an \(\ssee\)-removable ribbon in \(\tau\) by Lemma~\ref{xiuv}.
 
We have
\begin{align*}
\dist(u,w) = 
\dist(u,v) + \dist(v,z) + \dist(z,w) = 
 \dist(u,v) + \dist(z,w) + 1
\end{align*}
and
\begin{align*}
\res(z) = \res(\ea v) = \res(v) + \overline{1} = \res(u) + \overline{\dist(u,v) + 1}.
\end{align*}
It follows by (\ref{splithookcont}) and Lemma~\ref{hookcont} that 
\begin{align*}
\cont(\xi^\tau_{u,w}) &= \alpha(\res(u), \dist(u,w)+1) = \alpha(\res(u),  \dist(u,v) + \dist(z,w) + 2)\\
&= \alpha(\res(u), \dist(u,v)+1) + \alpha(\res(u) + \overline{\dist(u,v)+1}, \dist(z,w)+1)\\
&= \alpha(\res(u),  \dist(u,v)+1) + \alpha(\res(z), \dist(z,w)+1)\\
&= \cont(\xi_1) + \cont(\xi_2).
\end{align*}

Assume by way of contradiction that \(\cont(\xi_1) \succ \cont(\xi_2)\). Then we have \(\cont(\xi_1) \succ \cont(\xi^\tau_{u,w}) \succ \cont(\xi_2)\) by Lemma~\ref{GenConvexLem}. But, since \(\xi^\tau_{u,w}\) is \(\ssee\)-removable in \(\tau\), this contradicts \(\xi_1\) being a minimal \(\ssee\)-removable ribbon in \(\tau\). Therefore \(\cont(\xi_2) \succeq \cont(\xi_1)\), as desired.

 {\em Case 2: Assume \(w \nearrow u\).} Note that since \(\xi_1 \sqcup \xi_2\) is connected, we must have \(w \in \{\so u, \we u\}\). But, since \((u,v) \in \textup{Rem}_\tau\),
we have \(\so u \notin \tau\), so \(w=\we u\). We have by Lemma~\ref{remhookswap}\eqref{cond:remhookswap2} that \((z,v) \in \textup{Rem}_\tau\). Therefore \(\xi^\tau_{z,v}\) is an \(\ssee\)-removable ribbon in \(\tau\) by Lemma~\ref{xiuv}. 
 It can be shown along the lines of Case 1 that \(\cont(\xi^\tau_{z,v}) = \cont(\xi_1) + \cont(\xi_2)\). Thus if \(\cont(\xi_1) \succ \cont(\xi_2)\), we derive a contradiction along the same lines as Case 1 as well, so we have \(\cont(\xi_2) \succeq \cont(\xi_1)\), as desired.
 
 {\em Case 3: Assume \(u \nearrow z \nearrow v \nearrow w\)}. It follows from Lemma~\ref{remhookswap}\eqref{cond:remhookswap3} that \((\ssee z, w) \in \textup{Rem}_\tau\). 
Therefore \(\xi^\tau_{\ssee z,w}\) is an \(\ssee\)-removable ribbon in \(\tau\) by Lemma~\ref{xiuv}.
 We have
\(
 \dist(\ssee z ,w) = \dist(z,w),
\)
 and
\(
 \res(\ssee z) =\res(z),
\)
 so by Lemma~\ref{hookcont} it follows that 
 \begin{align*}
 \cont(\xi^\tau_{\ssee z,w}) &= \alpha(\res(\ssee z), \dist(\ssee z,w)+1)\\
 &= \alpha(\res(z),\dist(z,w)+1) = \cont(\xi_2).
 \end{align*}
As \(\xi_1\) is a minimal \(\ssee\)-removable ribbon in \(\tau\), we have then that
\(
\cont(\xi_2) = \cont(\xi^\tau_{\ssee z,w}) \succeq \cont(\xi_1),
\)
as desired.

{\em Case 4: Assume \(z \nearrow u \nearrow w \nearrow v\)}. 
It follows from Lemma~\ref{remhookswap}\eqref{cond:remhookswap4} that \((z, \ssee w) \in \textup{Rem}_\tau\).
Therefore \(\xi^\tau_{z, \ssee w}\) is an \(\ssee\)-removable ribbon in \(\tau\) by Lemma~\ref{xiuv}.
As in Case 3, it is straightforward to show that \(\cont(\xi^\tau_{z, \ssee w}) = \cont(\xi_2)\), and thus by the minimality of \(\xi_1\) we have
\(
\cont(\xi_2) = \cont(\xi^\tau_{z, \ssee w}) \succeq \cont(\xi_1),
\)
as desired.

{\em Case 5: Assume \(u \nearrow z \nearrow w \nearrow v\)}. It follows from Lemma~\ref{remhookswap}\eqref{cond:remhookswap5} that \((\ssee z, \ssee w) \in \textup{Rem}_\tau\tau\). Then \(\xi^\tau_{\ssee z, \ssee w}\) is an \(\ssee\)-removable ribbon in \(\tau\). 

We have
\(
 \dist(\ssee z,\ssee w) = \dist (z,w),
\)
 and
\(
 \res(\ssee z) = \res(z)
\), so it follows by Lemma~\ref{hookcont} that 
 \begin{align*}
 \cont(\xi^\tau_{\ssee z,\ssee w}) &= \alpha(\res(\ssee z), \dist(\ssee z,\ssee w)+1)\\
 &= \alpha(\res(z),\dist(z,w)+1) = \cont(\xi_2).
 \end{align*}
As \(\xi_1\) is a minimal \(\ssee\)-removable ribbon in \(\tau\), it follows that
\(
\cont(\xi_2) = \cont(\xi^\tau_{\ssee(z),w}) \succeq \cont(\xi_1),
\)
as desired.

{\em Case 6: Assume \(z \nearrow u \nearrow v \nearrow w\)}.  It follows from Lemma~\ref{remhookswap}\eqref{cond:remhookswap6} that \((z, w) \in \textup{Rem}_\tau\tau\). Thus \(\xi^\tau_{z,w}\) is an \(\ssee\)-removable ribbon in \(\tau\). 
But then by Lemma~\ref{xiuv} and the minimality of \(\xi_1\) we have
\(
\cont(\xi_2) =\cont(\xi^\tau_{z,w}) \succeq \cont(\xi_1),
\)
as desired, completing the proof. 
\end{proof}
